% ═══════════════════════════════════════════════════════════════
% ARBEITSBLATT: Das Ohmsche Gesetz (Klasse 8) - ZWEISPALTIG
% Differenziert: * (Basis) / ** (Standard) / *** (Erweitert)
% ═══════════════════════════════════════════════════════════════

\documentclass[10pt, a4paper]{article}

% ═══════════════════════════════════════════════════════════════
% PAKETE
% ═══════════════════════════════════════════════════════════════

\usepackage[utf8]{inputenc}
\usepackage[T1]{fontenc}
\usepackage[ngerman]{babel}
\usepackage{helvet}
\renewcommand{\familydefault}{\sfdefault}

\usepackage[margin=1.5cm, top=2cm, bottom=1.5cm]{geometry}
\usepackage{amsmath, amssymb}
\usepackage{siunitx}
\sisetup{locale = DE, per-mode = symbol}

\usepackage{graphicx}
\usepackage{tikz}
\usetikzlibrary{circuits.ee.IEC, positioning}
\usepackage{enumitem}
\usepackage{tabularx}
\usepackage{booktabs}
\usepackage{array}
\usepackage{multicol}

\usepackage{xcolor}
\usepackage{tcolorbox}
\usepackage{fancyhdr}
\usepackage{lastpage}

% Kompaktere Listen
\setlist{nosep, leftmargin=1.2em}

% Spaltenabstand
\setlength{\columnsep}{1cm}
\setlength{\columnseprule}{0.4pt}

% ═══════════════════════════════════════════════════════════════
% FARBEN
% ═══════════════════════════════════════════════════════════════

\definecolor{physik}{HTML}{2c5aa0}
\definecolor{grau}{HTML}{888888}
\definecolor{hellgrau}{HTML}{f0f0f0}
\definecolor{success}{HTML}{2a7a4b}
\definecolor{stern1}{HTML}{2a7a4b}
\definecolor{stern2}{HTML}{b35c00}
\definecolor{stern3}{HTML}{c41e3a}

% ═══════════════════════════════════════════════════════════════
% VARIABLEN
% ═══════════════════════════════════════════════════════════════

\newcommand{\thema}{Das Ohmsche Gesetz}
\newcommand{\fach}{Physik}
\newcommand{\klasse}{8}

% ═══════════════════════════════════════════════════════════════
% KOPF- UND FUSSZEILE
% ═══════════════════════════════════════════════════════════════

\pagestyle{fancy}
\fancyhf{}
\renewcommand{\headrulewidth}{0.4pt}
\renewcommand{\footrulewidth}{0pt}

\lhead{\small\textcolor{physik}{\textbf{\fach}} Kl.\,\klasse{} -- \thema}
\rhead{\small Name: \rule{4cm}{0.4pt}}

\cfoot{\small\thepage/\pageref{LastPage}}

% ═══════════════════════════════════════════════════════════════
% BEFEHLE
% ═══════════════════════════════════════════════════════════════

\newcommand{\luecke}[1]{\rule{#1}{0.4pt}}
\newcommand{\antwort}[1]{\vspace{#1}\hrule\vspace{2pt}}

\newtcolorbox{formelbox}{
    colback=hellgrau,
    colframe=physik,
    boxrule=1pt,
    arc=0pt,
    left=6pt, right=6pt, top=6pt, bottom=6pt
}

\newtcolorbox{merksatz}{
    colback=hellgrau,
    colframe=success,
    boxrule=1pt,
    arc=0pt,
    left=6pt, right=6pt, top=4pt, bottom=4pt,
    fonttitle=\bfseries,
    title=Merksatz
}

\newcommand{\niveauEins}{\textcolor{stern1}{\textbf{(*)}}}
\newcommand{\niveauZwei}{\textcolor{stern2}{\textbf{(**)}}}
\newcommand{\niveauDrei}{\textcolor{stern3}{\textbf{(***)}}}
\newcommand{\punkte}[1]{\hfill\textcolor{grau}{\small[#1 BE]}}

% ═══════════════════════════════════════════════════════════════
% DOKUMENT
% ═══════════════════════════════════════════════════════════════

\begin{document}

% ═══════════════════════════════════════════════════════════════
% KOPFBEREICH (einspaltig)
% ═══════════════════════════════════════════════════════════════

\begin{center}
    {\large\textbf{\textcolor{physik}{Arbeitsblatt: \thema}}}
\end{center}

\vspace{0.3em}

\begin{tcolorbox}[colback=hellgrau, colframe=grau, boxrule=0.5pt, arc=0pt, left=4pt, right=4pt, top=2pt, bottom=2pt]
\small
\niveauEins{} Basis \quad
\niveauZwei{} Standard \quad
\niveauDrei{} Erweitert
\hfill
\textbf{Punkte:} \luecke{0.8cm} / 25 BE
\end{tcolorbox}

\vspace{0.5em}

% ═══════════════════════════════════════════════════════════════
% FORMEL + MERKSATZ (einspaltig)
% ═══════════════════════════════════════════════════════════════

\begin{multicols}{2}

\section*{1. Formel und Merksatz}

\begin{formelbox}
\textbf{Das Ohmsche Gesetz:}
\begin{center}
{\large $U = R \cdot I$}
\end{center}
\vspace{0.3em}
\small
\begin{tabular}{lll}
$U$ & = & Spannung in \luecke{1.5cm} \\[0.4em]
$R$ & = & \luecke{2cm} in Ohm \\[0.4em]
$I$ & = & Stromstärke in \luecke{1.5cm} \\
\end{tabular}
\end{formelbox}

\vspace{0.3em}

\begin{merksatz}
Je \luecke{1.5cm} die Spannung, desto \luecke{1.5cm} die Stromstärke.
Der Widerstand bleibt \luecke{2cm}.
\end{merksatz}

\columnbreak

% ═══════════════════════════════════════════════════════════════
% MESSWERTE
% ═══════════════════════════════════════════════════════════════

\section*{2. Messwerte \niveauEins}

\textbf{Widerstand:} $R = \SI{20}{\ohm}$

\vspace{0.3em}

\small
\begin{tabular}{|c|c|c|c|c|c|}
\hline
$U$ in V & 2 & 4 & 6 & 8 & 10 \\
\hline
$I$ in A & & & & & \\[0.8em]
\hline
\end{tabular}
\normalsize

\vspace{0.5em}

\textbf{Zeichne die Messgerade:}

\begin{tikzpicture}[scale=0.45]
    \draw[step=0.5cm, gray!30, very thin] (0,0) grid (8,6);
    \draw[step=1cm, gray!50, thin] (0,0) grid (8,6);
    \draw[thick, ->] (0,0) -- (8.3,0) node[right] {\tiny$U$/V};
    \draw[thick, ->] (0,0) -- (0,6.3) node[above] {\tiny$I$/A};
    \foreach \x in {0,2,4,6,8,10} {
        \draw (\x/1.25,0) -- (\x/1.25,-0.1);
        \node[below] at (\x/1.25,-0.15) {\tiny\x};
    }
    \foreach \y/\label in {0/0, 2/0.2, 4/0.4} {
        \draw (0,\y) -- (-0.1,\y);
        \node[left] at (-0.15,\y) {\tiny\label};
    }
\end{tikzpicture}

\end{multicols}

% ═══════════════════════════════════════════════════════════════
% AUFGABEN (zweispaltig)
% ═══════════════════════════════════════════════════════════════

\begin{multicols}{2}

\section*{3. Aufgaben}

\subsection*{\niveauEins{} Basis}

\textbf{Aufgabe 1:} Lückentext \punkte{3}

Wenn man die Spannung \textbf{verdoppelt}, dann \luecke{2cm} sich auch die \luecke{2cm}.

Die Stromstärke ist \luecke{2cm} zur Spannung.

\vspace{0.8em}

\textbf{Aufgabe 2:} Berechne $U$. \punkte{2}

Gegeben: $R = \SI{10}{\ohm}$, $I = \SI{2}{\ampere}$

$U = R \cdot I = $ \luecke{1cm} $\cdot$ \luecke{1cm} $=$ \luecke{1.5cm}

\vspace{0.8em}

\subsection*{\niveauZwei{} Standard}

\textbf{Aufgabe 3:} Berechne. \punkte{6}

a) $R = \SI{50}{\ohm}$, $I = \SI{0,2}{\ampere}$

\quad $U = $ \luecke{2.5cm}

\vspace{0.5em}

b) $U = \SI{12}{\volt}$, $R = \SI{4}{\ohm}$

\quad $I = U/R = $ \luecke{2.5cm}

\vspace{0.5em}

c) $U = \SI{10}{\volt}$, $I = \SI{0,5}{\ampere}$

\quad $R = U/I = $ \luecke{2.5cm}

\vspace{0.8em}

\textbf{Aufgabe 4:} Ordne zu. \punkte{3}

\begin{tabular}{lcl}
$U$ & \luecke{1.2cm} & Ampere \\
$I$ & \luecke{1.2cm} & Ohm \\
$R$ & \luecke{1.2cm} & Volt \\
\end{tabular}

\columnbreak

\subsection*{\niveauDrei{} Erweitert}

\textbf{Aufgabe 5:} Fehlersuche \punkte{4}

Ein Schüler rechnet:

``$R = \SI{50}{\ohm}$, $I = \SI{200}{\milli\ampere}$.

$U = 50 \cdot 200 = \SI{10000}{\volt}$.''

a) Fehler?

\antwort{0.8cm}
\antwort{0.8cm}

b) Richtig:

\vspace{1.2cm}

\textbf{Aufgabe 6:} Transfer \punkte{4}

$R = \SI{20}{\ohm}$, $U = \SI{6}{\volt}$

a) Berechne $I$:

\vspace{1.2cm}

b) $U$ verdoppelt sich auf \SI{12}{\volt}. Was passiert mit $I$? Begründe.

\antwort{0.8cm}
\antwort{0.8cm}

\vspace{0.5em}

\textbf{Aufgabe 7:} Begründung \punkte{3}

Warum bleibt $R$ konstant, auch wenn $U$ sich ändert?

\antwort{0.8cm}
\antwort{0.8cm}
\antwort{0.8cm}

\end{multicols}

% ═══════════════════════════════════════════════════════════════
% PUNKTEÜBERSICHT (einspaltig)
% ═══════════════════════════════════════════════════════════════

\vspace{0.5em}

\begin{tcolorbox}[colback=hellgrau, colframe=grau, boxrule=0.5pt, arc=0pt, left=4pt, right=4pt, top=2pt, bottom=2pt]
\small
\textbf{Punkte:}
\niveauEins{} 1--2 = 5 BE \quad
\niveauZwei{} 3--4 = 9 BE \quad
\niveauDrei{} 5--7 = 11 BE \quad
\textbf{Gesamt: 25 BE}
\hfill
\textbf{Erreicht:} \luecke{1cm} BE
\end{tcolorbox}

% ═══════════════════════════════════════════════════════════════
% SEITE 2: FORMELKARTE (einspaltig)
% ═══════════════════════════════════════════════════════════════

\newpage

\begin{center}
    {\large\textbf{\textcolor{physik}{Formelkarte: Ohmsches Gesetz}}}
\end{center}

\vspace{0.5em}

\begin{multicols}{2}

\begin{formelbox}
\begin{center}
\textbf{Die drei Formeln:}

\vspace{0.5em}

\begin{tabular}{|c|c|c|}
\hline
$U = R \cdot I$ & $I = \dfrac{U}{R}$ & $R = \dfrac{U}{I}$ \\[0.8em]
\hline
\end{tabular}

\vspace{0.8em}

\textbf{Einheiten:}

\begin{tabular}{ll}
$U$ in Volt & (V) \\
$I$ in Ampere & (A) \\
$R$ in Ohm & ($\Omega$) \\
\end{tabular}

\vspace{0.5em}

$\SI{1000}{\milli\ampere} = \SI{1}{\ampere}$

\end{center}
\end{formelbox}

\end{multicols}

% URI-Dreieck (einspaltig, volle Breite)
\begin{center}
\begin{tikzpicture}[scale=4.0]
    % Dreieck
    \draw[very thick, physik] (0,0) -- (4,0) -- (2,3) -- cycle;
    % Horizontale Trennlinie (U oben, R|I unten)
    \draw[very thick, physik] (0.67,1) -- (3.33,1);
    % Vertikale Trennlinie (R links, I rechts)
    \draw[very thick, physik] (2,0) -- (2,1);
    % Beschriftungen
    \node at (2, 1.9) {\LARGE\textbf{$U$}};
    \node at (1, 0.4) {\LARGE\textbf{$R$}};
    \node at (3, 0.4) {\LARGE\textbf{$I$}};
\end{tikzpicture}

\vspace{0.3em}
\textbf{URI-Dreieck:} Gesuchtes zudecken $\rightarrow$ Rest ablesen
\end{center}

\vspace{0.5em}

\begin{multicols}{2}

\begin{tcolorbox}[colback=white, colframe=physik, boxrule=1pt, arc=0pt, title=\textbf{Beispiel}]

\textbf{Gegeben:} $U = \SI{12}{\volt}$, $R = \SI{4}{\ohm}$

\textbf{Gesucht:} $I = \,?$

\textbf{Formel:} $I = \dfrac{U}{R}$

\textbf{Einsetzen:} $I = \dfrac{\SI{12}{\volt}}{\SI{4}{\ohm}} = \SI{3}{\ampere}$

\textbf{Antwort:} Die Stromstärke beträgt $\SI{3}{\ampere}$.

\end{tcolorbox}

\vspace{1em}

\begin{tcolorbox}[colback=white, colframe=stern2, boxrule=1pt, arc=0pt, title=\textbf{Typische Fehler}]
\small
\begin{itemize}[leftmargin=1em]
    \item mA nicht in A umgerechnet
    \item Formel falsch umgestellt
    \item Einheiten vergessen
\end{itemize}
\end{tcolorbox}

\end{multicols}

\end{document}
